\pg{333} existencia e para a não-existencia. --- No ente \textbf{\emph{necessario,}} a \emph{existencia} é um elemento da \emph{essencia}, e identifica-se com esta; no ente \emph{contingente,} a \emph{existencia} é \emph{indifferente} para a \textbf{\emph{essencia,}} e uma coisa é \emph{realmente} distincta de outra \footnote{\emph{Ente necessario} é o que não pode deixar de ser. Na verdade, \emph{ente} é o que possue o \emph{ser}, é o que tem a \emph{existencia}, é o \emph{existente}. Portanto, o ente \emph{necessario} é o que existe \emph{necessariamente}, e por isso possue o \emph{ser} por tal modo, que lhe repugna o \emph{não-ser}. --- O \emph{ente contingente}, por ser \emph{contingente}, existe por tal modo, que pode deixar de existir, sem que d'ahi resulte algum absurdo ou inconveniente. --- O ente \emph{necessario} pode ser \emph{absoluto} e \emph{hypothetico}. É \emph{absoluto,} quando o \emph{não-ser} lhe repugna pela sua intrinseca natureza, independentemente de toda a condição ou hypothese; tal é \textsc{Deus}. É \emph{hypothetico}, quando o \emph{não-ser} lhe repugna, não pela sua intrinseca natureza, mas debaixo de uma certa condição ou hypothese; tal é toda a creatura, emquanto existe. --- Aqui fallamos do ente \emph{absoluto.}\par O ente \emph{absoluto,} não tendo nenhuma potencia ou indifferença para a existencia, mas existindo em virtude da propria essencia, é ente \emph{por si, pela sua essencia,} de modo que a sua essencia é a razão sufficiente da sua existencia. --- Note-se que quando se diz que o ente \emph{absoluto} existe pela \emph{sua essencia,} não se entende que a essencia d'esse ente é a \emph{causa} da sua existencia: o que é absurdo; mas deve entender-se que esse ente não depende, na sua existencia, de uma causa extrinseca e superior. --- Advirta-se tambem que, embora a \emph{essencia} do ente \emph{absoluto} se identifique com a \emph{existencia}, todavia a nossa intelligencia não percebe essas duas coisas por uma só idea, mas por duas ideas distinctas, uma das quaes representa a \emph{essencia}, e a outra a \emph{existencia}. E nada mais natural. A intelligencia humana não tem um conhecimento intuitivo do Ente \emph{absoluto}, mas eleva-se a Elle por meio das coisas creadas, e nestas a \emph{essencia é realmente} distincta da \emph{existencia}.}.
